\documentclass[11pt,letterpaper]{article}
\usepackage[T1]{fontenc}
\usepackage[utf8]{inputenc}
\usepackage[scaled=0.95]{helvet}
\renewcommand{\familydefault}{\sfdefault}
\usepackage[top=1.8cm,bottom=1.6cm,left=2.2cm,right=2.2cm]{geometry}
\usepackage{microtype}
\usepackage[hidelinks]{hyperref}
\pagestyle{empty}
\setlength{\parindent}{0pt}
\setlength{\parskip}{8pt}

\begin{document}

{\large\bfseries Grover Aruquipa}\par\vspace{-6pt}
{\small Qu\'ebec, QC, Canada \ \textbullet\ \ grover.aruquipa.9@gmail.com \ \textbullet\ \ groveraruquipa.github.io}

\vspace{6pt}
Hiring Team, Horizon Surgical Systems\\
Senior/Staff Mechanical Engineer, Robotic Hardware

Dear Hiring Team,

Most of my work so far has been about one problem: making a mechanism move precisely in a small space while it keeps control of what it holds. When I read about Polaris and the first cataract surgery performed with it, I could not think of a place where that problem matters more than inside a human eye. That is why I am applying, and why this position means a great deal to me.

Over the last three years at Universit\'e Laval, with Prof.\ Cl\'ement Gosselin and Prof.\ Patrice Lambert, I designed and built three robots whose kinematic structures are not standard arms: kinematically redundant parallel and hybrid mechanisms with 4, 7 and 8 actuators, all mounted on the base, with the gripper driven through the mechanism itself. I took each one myself from the kinematic model and the SolidWorks design through joint and transmission design, component selection, machining, printing, wiring and assembly to real-time control and testing. When something did not work, I was the one taking it apart on the bench.

The one I would most like to show you is a 4-DoF coaxial wrist-gripper. It passes four actuation channels through concentric hollow shafts, rolls its tool without limit while a non-backdrivable gripper keeps hold, and has 0.25\,kg of moving mass. On the bench it completed ten continuous revolutions with sub-degree tracking. Unlimited roll about the tool axis, with the jaw actuated from behind through the shaft, is close to what a surgical instrument has to do, and it is the kind of subsystem I would like to build at Horizon. Before Laval, my master's research at FEMTO-ST in France was on parallel robots for micro-manipulation, so working at very small scales is not new to me.

I want to be direct about what I still have to learn. I have not designed under an FDA quality system or for IEC 60601, and my years are in research rather than product development. I document every design for peer review, I write my own analysis in MATLAB, Python and C++, and I expect to learn your design controls quickly from the team. Since the posting notes that PhD work on robot mechanism design may substitute for industry experience, I would be glad to be considered at the Senior or Staff level, whichever your team finds fits best.

I defend my PhD in January 2027 and would be glad to work on site in Santa Monica or Malibu, five days a week, after that. I would need U.S. visa sponsorship and would appreciate knowing early whether that is possible. Prof.\ Gosselin, Prof.\ Lambert and Univ.-Prof.\ Andreas M\"uller of JKU Linz, who hosted my research stay this year, can speak to my work; their contact details are in my CV.

Thank you for your time. I would be happy to walk you through any of the prototypes, on video or in person.

Sincerely,\\
Grover Aruquipa

\end{document}
